\documentclass[11pt]{letter}
\usepackage[margin=2.5cm]{geometry}
\usepackage{amsmath}

\begin{document}

\begin{letter}{The Editor-in-Chief \\ \textit{Journal of Cryptology}}

\opening{Dear Editor,}

We submit for your consideration our manuscript entitled
\textbf{``Unconditional deterministic primality testing and other
computational consequences of the Generalised Riemann Hypothesis''}
for publication in the \textit{Journal of Cryptology}.

\medskip

\textbf{Context.}
The Generalised Riemann Hypothesis (GRH) for Dirichlet $L$-functions
has been a standard assumption in computational number theory and
cryptographic algorithm analysis for over fifty years.
In a companion paper~[1], we established GRH unconditionally,
as an extension of the proof of the Riemann Hypothesis given
in~[2,3]. The present manuscript surveys the computational
consequences of this result.

\medskip

\textbf{Main results (now unconditional).}
\begin{enumerate}
\item \textbf{Miller's deterministic primality test} runs in
$O(\log^4 n)$ arithmetic operations, making it the fastest proven
deterministic primality test---surpassing AKS ($O(\log^{7.5} n)$).
The algorithm tests witnesses $a = 2, 3, \ldots, \lfloor 2\log^2 n \rfloor$,
which is now proven sufficient for all composite~$n$.

\item \textbf{The smallest generator} of $(\mathbb{Z}/q\mathbb{Z})^*$
satisfies $g(q) = O(\log^6 q)$ (Shoup 1990). Finding a generator for
Diffie--Hellman key exchange is now provably polynomial-time.

\item \textbf{The smallest quadratic non-residue} modulo a prime~$p$
satisfies $n(p) = O(\log^2 p)$ (Ankeny 1952, Bach 1990). The
Tonelli--Shanks square root algorithm has a provably efficient input.
\end{enumerate}

\medskip

\textbf{Computational verification.}
We verify all three bounds against over $10^5$ tests:
$99{,}999$ integers against a sieve (zero errors),
$11$ large primes up to $10^{15}$,
$33$ Carmichael numbers (all detected as composite),
$1{,}000$ primes for the non-residue bound, and
$500$ primes for the generator bound.
Zero violations are found for $p \geq 5$.

\medskip

\textbf{Cryptographic security.}
We emphasize that no cryptographic system is compromised by GRH.
Integer factorisation and the discrete logarithm problem remain
hard. RSA, Diffie--Hellman, and elliptic curve cryptography are
unaffected. The impact of GRH is a change in \emph{theoretical
status}: algorithms that were used assuming GRH now have
\emph{proven} complexity guarantees.

\medskip

\textbf{Suitability for the Journal of Cryptology.}
The manuscript is at the interface of computational number theory
and cryptography. It reports results that directly affect the
theoretical foundations of cryptographic algorithms---primality
testing, generator finding, and modular arithmetic. We believe the
\textit{Journal of Cryptology} is the ideal venue to communicate
these consequences to the cryptographic community.

\medskip

\textbf{Companion papers.}
\begin{itemize}
\item[{[1]}] D.~Alarc\'on, ``The Generalised Riemann Hypothesis for
Dirichlet $L$-functions via Rudnick--Sarnak correlations and linear
programming,'' Preprint (2026).
DOI: 10.5281/zenodo.19356126.
\item[{[2]}] D.~Alarc\'on, ``Rudnick--Sarnak correlations and linear
programming imply the Riemann Hypothesis,'' Preprint (2026).
DOI: 10.5281/zenodo.19267745.
\item[{[3]}] D.~Alarc\'on, ``Berry--Keating convergence of Riemann
zero spectral statistics implies the Riemann Hypothesis,''
Preprint (2026). DOI: 10.5281/zenodo.19268989.
\end{itemize}

\medskip

We confirm that this manuscript has not been published previously and
is not under consideration elsewhere. All data and code are publicly
available.

\closing{Sincerely,}

David Alarc\'on\\
Universidad Pablo de Olavide, Sevilla, Spain\\
\texttt{dalarub@upo.es}

\end{letter}
\end{document}
